本文建立了从工业历史轨迹到世界模型、再到策略优化和 NeoRL 仿真验证的完整流程。实验比较了 MLP、GRU、LSTM、Transformer-2L 和 Transformer-4L，并在三个数据规模上统一评价单步与多步递归预测。综合预测结果选择 M100 的 GRU，以及 M1000、M10000 的 Transformer-2L 作为对应数据规模的冻结世界模型。

基于这些世界模型，本文完成 CEM、iCEM、MPPI 三种规划方法和 MB-PPO 的策略优化。最终 M100 使用 GRU+CEM，M1000 与 M10000 使用 Transformer-2L+MPPI；三者的1000步累计奖励相对 BC 分别提高14,229、68,282和64,535，表明世界模型能够在当前基准中有效支持长期策略优化。

实验还揭示了多步误差累积、预测精度与控制效果不完全一致、世界模型误差可能被优化器放大，以及数据增加不一定带来单调控制收益等问题。这些现象提示，真实工业应用除了追求预测精度和控制回报，还必须重视数据覆盖、不确定性估计、OOD 检测、安全约束、回退控制和持续监测，从而在可验证的安全边界内发挥世界模型的预测与决策价值。
